\chapter{Solutions to Chapter 8: Unification of the Fisherian and Neymanian Inferences in Randomized Experiments}
\label{sol:chapter08}

\begin{exercisesolution}{Re-analyzing \citet{angrist2009effects}'s data}{MPE 中 studentized FRT 的可复现框架}
本题是 Problem \ref{para::AL-mpe-neyman} 的 Fisherian 对应版本。实验单位是 schools，设计是 matched-pairs experiment；删去 pair 6 以及 noncompliance pairs 后，分析对象是 Chapter \ref{chapter::mpe} 中给出的 $14$ 个 complete pairs。对每个 outcome year，例如 \ri{pr01} 或 \ri{pr02}，先构造 pair-level treated-control difference
\[
d^Y_i=Y_{i,T}^{\mathrm{obs}}-Y_{i,C}^{\mathrm{obs}},
\]
并对 pretreatment covariates 构造同方向的 differences
\[
d^{99}_i=X^{99}_{i,T}-X^{99}_{i,C},\qquad
d^{00}_i=X^{00}_{i,T}-X^{00}_{i,C}.
\]
在 sharp null hypothesis $\HF$ 下，每个 pair 内两所学校的 observed outcomes 固定，randomization 只会独立翻转 pair-level difference 的符号。因此若 $s_i\in\{-1,1\}$ 表示相对于 observed assignment 是否翻转第 $i$ 个 pair，则 randomization distribution 由
\[
d^Y_i(s)=s_id^Y_i,\qquad
d^{99}_i(s)=s_id^{99}_i,\qquad
d^{00}_i(s)=s_id^{00}_i
\]
生成。这里 covariate differences 也随 sign 翻转，因为在每个 hypothetical assignment 下 treated-control 的方向随 assignment 改变。

不调整 covariates 时，推荐的 studentized statistic 是 intercept-only pair-level OLS 的 $t$-statistic：
\[
d^Y_i(s)=\alpha+u_i,\qquad
T_0(s)=\frac{\hat\alpha(s)}{\widehat{\se}\{\hat\alpha(s)\}}.
\]
由于 $\hat\alpha(s)=n^{-1}\sum_i d^Y_i(s)$ 且 usual OLS intercept standard error 等于
\[
\left\{\frac{1}{n(n-1)}\sum_{i=1}^n
\{d^Y_i(s)-\bar d^Y(s)\}^2\right\}^{1/2},
\]
这正是 MPE 的 paired $t$ statistic。它和 \cref{prop::mpe-variance,thm::varianceest-mp} 中的 Neymanian variance estimator 使用同一份 pair-level variation。

调整 covariates 时，按 \cref{prop::colin2018regression} 拟合
\[
d^Y_i(s)=\alpha+\gamma_1d^{99}_i(s)+\gamma_2d^{00}_i(s)+u_i,
\]
并使用 intercept 的 $t$-statistic
\[
T_{\mathrm{adj}}(s)=
\frac{\hat\alpha(s)}{\widehat{\se}\{\hat\alpha(s)\}}.
\]
这对应 Chapter \ref{chapter::unification-fisher-neyman} 中对 MPE 的建议：没有 covariates 时使用 intercept-only regression 的 studentized statistic；有 covariates 时使用 pair-level covariate-adjusted regression 的 intercept $t$-statistic。其理由与 \cref{sec::strong-weak-cre-frt} 平行：studentization 让 FRT 在 $\HF$ 下保留 finite-sample exactness，并在 weak null $\HN$ 下具有更好的 large-sample 行为；covariate adjustment 的角色则是利用 pretreatment Bagrut passing rates 提高 power。

因为本题要求报告 $p_\frt$，但这里不伪造没有实际运行数据得到的数值。可复现的 \ri{R} 框架如下；如果使用教材中的 \ri{AL2009.csv}，它会同时给出 \ri{pr01} 和 \ri{pr02}、未调整和调整后的 two-sided $p_\frt$。对于 \ri{pr02} 中 pair 13 的缺失值，主分析采用 complete-pair rule：删除整个 pair，而不是只删除一个 school。这样做尊重 MPE 的 randomization unit。
\begin{lstlisting}
make_pair_data <- function(dat, yname) {
  pieces <- by(dat, dat$pair, function(dd) {
    dd <- dd[order(dd$z), ]
    y0 <- dd[dd$z == 0, yname]
    y1 <- dd[dd$z == 1, yname]
    x990 <- dd[dd$z == 0, "pr99"]
    x991 <- dd[dd$z == 1, "pr99"]
    x000 <- dd[dd$z == 0, "pr00"]
    x001 <- dd[dd$z == 1, "pr00"]
    c(dy = y1 - y0, dx99 = x991 - x990, dx00 = x001 - x000)
  })
  out <- as.data.frame(do.call(rbind, pieces))
  na.omit(out)
}

t_unadj <- function(D) {
  summary(lm(dy ~ 1, data = D))$coef[1, "t value"]
}

t_adj <- function(D) {
  summary(lm(dy ~ dx99 + dx00, data = D))$coef[1, "t value"]
}

enumerate_signs <- function(n) {
  as.matrix(expand.grid(rep(list(c(-1, 1)), n)))
}

frt_mpe <- function(D, stat) {
  D <- na.omit(D)
  obs <- stat(D)
  signs <- enumerate_signs(nrow(D))
  vals <- apply(signs, 1, function(s) {
    Ds <- D
    Ds$dy <- s * D$dy
    Ds$dx99 <- s * D$dx99
    Ds$dx00 <- s * D$dx00
    stat(Ds)
  })
  c(t.obs = obs,
    p.frt = mean(abs(vals) >= abs(obs)),
    n.pairs = nrow(D))
}

dat <- read.csv("AL2009.csv")
D01 <- make_pair_data(dat, "pr01")
D02 <- make_pair_data(dat, "pr02")

frt_mpe(D01, t_unadj)
frt_mpe(D01, t_adj)
frt_mpe(D02, t_unadj)
frt_mpe(D02, t_adj)
\end{lstlisting}

如果 pairs 很少，可以像上面那样枚举所有 $2^n$ 个 assignments，得到 exact $p_\frt$。如果 pairs 更多，则把 \ri{enumerate\_signs} 换成 Monte Carlo signs，并使用
\[
\tilde p_\frt=
\frac{1+\sum_{r=1}^R \mathbb{I}\{|T_r|\geq |T_{\mathrm{obs}}|\}}{1+R}
\]
作为 finite-sample valid 的 Monte Carlo approximation；其逻辑见 Chapter \ref{ch::frt-cre} 中 Monte Carlo FRT 的讨论。

\emph{Reference / 用途：} 本题连接了 \cref{para::AL-mpe-neyman} 的 Neymanian analysis 与本章的 Fisherian recommendation。实践报告应同时列出 outcome year、是否调整 covariates、使用的 pair 数、observed studentized statistic 和 $p_\frt$。若 covariate adjustment 后 $p_\frt$ 明显变小，应解释为 pretreatment passing rates 对 outcome 有预测力，从而提高了 test statistic 的 sensitivity；不能解释为 regression model 自身保证了 validity。Validity 仍来自 MPE 的 randomization distribution。
\end{exercisesolution}

\begin{exercisesolution}{Replication of \citet{zhao2020caFRT}'s Figure 1}{SRE 中不同 FRT statistics 的 size simulation}
\citet{zhao2020caFRT} 的 Figure 1 比较了 SRE 下不同 FRT statistics 的有限样本表现。复现的核心不是某个单一数据集，而是一个 repeated-simulation workflow：先固定一个 SRE design 和一组 potential outcomes，再反复抽取 realized assignments；在每个 realized dataset 上计算若干 FRT $p$-values；最后在 weak null $\HN:\tau=0$ 下比较 empirical rejection rates。若某个 test 在 nominal level $\alpha$ 下的 rejection rate 大于 $\alpha$，它在该 data-generating process 下反保守；若不超过 $\alpha$，则 finite-sample 或 large-sample 上更稳妥。

对第 $k$ 个 stratum，记其大小为 $n_{[k]}$，treatment 数为 $n_{[k]1}$，control 数为 $n_{[k]0}$，权重为 $\pi_{[k]}=n_{[k]}/n$。未调整的 SRE estimator 和 conservative variance estimator 为
\[
\hat\tau_\textsc{S}
=\sum_{k=1}^K\pi_{[k]}\hat\tau_{[k]},
\]
\[
\hat V_\textsc{S}
=\sum_{k=1}^K\pi_{[k]}^2
\left\{
\frac{\hat S_{[k]}^2(1)}{n_{[k]1}}
+\frac{\hat S_{[k]}^2(0)}{n_{[k]0}}
\right\},
\]
见 \cref{sec::sre-neyman}。相应 studentized statistic 是
\[
t_\textsc{S}=\frac{\hat\tau_\textsc{S}}{\sqrt{\hat V_\textsc{S}}}.
\]
有额外 covariates 时，在每个 stratum 内使用 \citet{lin2013} regression adjustment，再按权重汇总：
\[
\hat\tau_{\textsc{L},\textsc{S}}
=\sum_{k=1}^K\pi_{[k]}\hat\tau_{\textsc{L},[k]},
\qquad
\hat V_{\textsc{L},\textsc{S}}
=\sum_{k=1}^K\pi_{[k]}^2\hat V_{\textsc{ehw},[k]},
\]
其中 stratum 内 covariates 必须按该 stratum 自身均值中心化；见 \cref{sec::covariate-SRE}。推荐 statistic 为
\[
t_{\textsc{L},\textsc{S}}
=\frac{\hat\tau_{\textsc{L},\textsc{S}}}
{\sqrt{\hat V_{\textsc{L},\textsc{S}}}}.
\]
本章 \cref{sec::xfrt-cre} 和 General recommendations 的结论说明，studentized statistics 尤其是 $t_{\textsc{L},\textsc{S}}$ 的优势在于同时兼顾 strong null 下的 randomization validity、weak null 下的 asymptotic validity，以及 covariates 有预测力时的 power。

下面给出复现 Figure 1 的完整代码骨架。它故意把 data-generating process 写成可替换的函数；实际复现时应把 \ri{make\_finite\_population} 中的参数改成 \citet{zhao2020caFRT} Figure 1 使用的设定，并设置相同的 \ri{B}、\ri{R}、nominal levels 和 random seed。
\begin{lstlisting}
library(sandwich)

make_finite_population <- function(K = 5, m = 20) {
  n <- K * m
  strata <- rep(seq_len(K), each = m)
  x <- rnorm(n) + rep(seq(-1, 1, length.out = K), each = m)
  y0 <- 1 + strata / K + x + rnorm(n)
  tau_i <- 0.5 * (x - mean(x))
  y1 <- y0 + tau_i - mean(tau_i)
  data.frame(strata = strata, x = x, y0 = y0, y1 = y1)
}

draw_sre <- function(dat, n1_each) {
  z <- integer(nrow(dat))
  for (k in sort(unique(dat$strata))) {
    id <- which(dat$strata == k)
    z[sample(id, n1_each[k])] <- 1
  }
  z
}

tau_s <- function(y, z, strata) {
  n <- length(y)
  out <- by(seq_along(y), strata, function(id) {
    nk <- length(id)
    pi_k <- nk / n
    pi_k * (mean(y[id][z[id] == 1]) - mean(y[id][z[id] == 0]))
  })
  sum(unlist(out))
}

v_s <- function(y, z, strata) {
  n <- length(y)
  out <- by(seq_along(y), strata, function(id) {
    nk <- length(id)
    pi_k <- nk / n
    yk <- y[id]
    zk <- z[id]
    pi_k^2 * (var(yk[zk == 1]) / sum(zk == 1) +
              var(yk[zk == 0]) / sum(zk == 0))
  })
  sum(unlist(out))
}

t_s <- function(y, z, strata) {
  tau_s(y, z, strata) / sqrt(v_s(y, z, strata))
}

lin_one_stratum <- function(y, z, x) {
  xc <- x - mean(x)
  fit <- lm(y ~ z * xc)
  b <- coef(fit)["z"]
  v <- vcovHC(fit, type = "HC2")["z", "z"]
  c(est = b, var = v)
}

t_ls <- function(y, z, x, strata) {
  n <- length(y)
  ans <- by(seq_along(y), strata, function(id) {
    nk <- length(id)
    pi_k <- nk / n
    tmp <- lin_one_stratum(y[id], z[id], x[id])
    c(est = pi_k * tmp["est"], var = pi_k^2 * tmp["var"])
  })
  mat <- do.call(rbind, ans)
  sum(mat[, "est"]) / sqrt(sum(mat[, "var"]))
}

permute_sre <- function(z, strata) {
  zp <- z
  for (k in sort(unique(strata))) {
    id <- which(strata == k)
    zp[id] <- sample(z[id])
  }
  zp
}

p_frt <- function(y, z, x, strata, stat, R = 2000) {
  obs <- stat(y, z, x, strata)
  vals <- replicate(R, {
    zp <- permute_sre(z, strata)
    stat(y, zp, x, strata)
  })
  (1 + sum(abs(vals) >= abs(obs))) / (R + 1)
}

stat_tau_s <- function(y, z, x, strata) tau_s(y, z, strata)
stat_t_s <- function(y, z, x, strata) t_s(y, z, strata)
stat_t_ls <- function(y, z, x, strata) t_ls(y, z, x, strata)

one_run <- function(pop, n1_each, R = 2000) {
  z <- draw_sre(pop, n1_each)
  y <- ifelse(z == 1, pop$y1, pop$y0)
  c(tauS = p_frt(y, z, pop$x, pop$strata, stat_tau_s, R),
    tS = p_frt(y, z, pop$x, pop$strata, stat_t_s, R),
    tLS = p_frt(y, z, pop$x, pop$strata, stat_t_ls, R))
}

set.seed(20260208)
pop <- make_finite_population(K = 5, m = 20)
n1_each <- tapply(pop$strata, pop$strata, length) / 2

B <- 1000
R <- 2000
pvals <- replicate(B, one_run(pop, n1_each, R))
levels <- c(0.01, 0.05, 0.10)
size <- sapply(levels, function(a) rowMeans(pvals <= a))
matplot(levels, t(size), type = "b", pch = 1:3, lty = 1,
        xlab = "nominal level", ylab = "empirical rejection rate")
abline(0, 1, lty = 2)
legend("topleft", legend = rownames(size), pch = 1:3, lty = 1)
\end{lstlisting}

上面的代码包括三类 FRT：以 $\hat\tau_\textsc{S}$ 为 statistic 的未 studentized FRT、以 $t_\textsc{S}$ 为 statistic 的 studentized FRT、以及以 $t_{\textsc{L},\textsc{S}}$ 为 statistic 的 covariate-adjusted studentized FRT。若要更贴近 \citet{zhao2020caFRT} 的 Figure 1，还应加入他们图中列出的其他 statistics，例如未 studentized 的 covariate-adjusted estimator、pooled regression coefficient 或没有正确分层的 permutation statistic；实现方式是在 \ri{p\_frt} 中替换 \ri{stat} 函数即可。

读图时应关注三点。第一，在 strong null 下，只要 statistic 对 observed assignment 和 permuted assignments 使用同一规则，FRT 的 $p_\frt$ 有 randomization validity；这是 Chapter \ref{ch::frt-cre} 的基本结论。第二，在 weak null 但存在 effect heterogeneity 时，未 studentized statistics 的 permutation distribution 可能与真实 sampling distribution 不匹配，机制见 \cref{sec::strong-weak-cre-frt}。第三，$t_{\textsc{L},\textsc{S}}$ 同时使用 SRE 的分层结构、covariate adjustment 和 studentization，因此是 \citet{zhao2020caFRT} 推荐的分析工具。

\emph{Remark / 用途：} 这道题的实务价值在于把“复现一张图”转化为可审计的 simulation protocol：写清 finite population、assignment mechanism、statistics、FRT 计算方式、Monte Carlo 次数和 rejection-rate 汇总方式。只要这些元素固定，图中的每个点都可复现；若没有原文 Figure 1 的精确 DGP，就只能报告复现框架和对不同 statistics 的理论预期，不能编造 empirical size 数值。
\end{exercisesolution}
